For $\Delta\in\Z$ and a prime $p$ define $\chi_\Delta(p)\in\{-1,0,1\}$ as follows. If $p$ is odd, $\chi_\Delta(p)$ is the Legendre symbol
$\bigl(\frac\Delta p\bigr)$: it is $0$ if $p\mid\Delta$, $1$ if $\Delta$ is a nonzero square modulo $p$, and $-1$ otherwise. If $p=2$, then
$\chi_\Delta(2)=0$ if $\Delta$ is even, $\chi_\Delta(2)=1$ if $\Delta\equiv\pm1\pmod 8$, and $\chi_\Delta(2)=-1$ if $\Delta\equiv\pm3\pmod 8$.
